\begin{python0}
from solutions import *; clear()
\end{python0}
\sectionthree{``When should I inline a function?''}

Now you might say: ``In that case, let's inline every function! That will make my program super fast!''

Well ...

Yes, you can inline any function.

The problem is that the your program might become larger than what it should be.

Good choices for functions to be changed to inline functions are those which are very small, especially those which are used frequently. Function bodies with only one or two lines of code are good candidates.

Now for some caveats \ldots{}

\begin{itemize}
\item If the compiler feels that the function body is too huge, it might
  actually ignore your request to inline a function. So an inline
  function is only a \textbf{\textit{request}} to the compiler to inline
  the function.
\item Some C++ compilers actually honor inline requests unless told it to do
  otherwise.
\item If too many functions are inlined, the executable binary code can be
  huge, which of course means that you need more memory to run the code
  which means that you might be slowing it down!!! If the executable
  binary code is too huge it might cause memory thrashing where memory
  pages constantly gets swapped in and out of your computer memory.
\end{itemize}


